\chapter[Chapter \thechapter]{}
\lettrine{M}{r} Crofts, excusably enough, said, <I told you so>; Sir~Impey Biggs observed curtly, <Very unfortunate.>

\zz
To chronicle Lord~Peter Wimsey's daily life during the ensuing week would be neither kind nor edifying. An enforced inactivity will produce irritable symptoms in the best of men. Nor did the imbecile happiness of Chief-Inspector Parker and Lady Mary Wimsey tend to soothe him, accompanied as it was by tedious demonstrations of affection for himself. Like the man in Max Beerbohm's story, Wimsey <hated to be touching.> He was only moderately cheered by hearing from the industrious Freddy Arbuthnot that Mr~Norman Urquhart was found to be more or less deeply involved in the disasters of the Megatherium Trust.

Miss~Kitty Climpson, on the other hand, was living in what she herself liked to call a <whirl of activity.> A letter, written the second day after her arrival in Windle, furnishes us with a wealth of particulars.

\begin{mail}{Hillside View,\\Windle, Westmorland.\\1st Jan. 1930.}{My dear Lord~Peter,}

I feel sure you will be anxious to hear, at the \emph{earliest possible} moment \emph{how} things are \emph{going}, and though I have only been here \emph{one} day, I really think I have \emph{not} done so \emph{badly}, all things considered!

My train got in quite late on Monday night, after a \emph{most dreary} journey, with a \emph{lugubrious} wait at \emph{Preston}, though thanks to your kindness in insisting that I should travel \emph{First-class}, I was not really at all tired! Nobody can realise what a \emph{great} difference these extra comforts make, especially when one is \emph{getting on} in years, and after the \emph{uncomfortable} travelling which I had to endure in my days of poverty, I feel that I am living in almost \emph{sinful} luxury! The carriage was \emph{well} heated—indeed, \emph{too much} so and I should have liked the window down, but that there was a \emph{very fat} business man, \emph{muffled} up to the eyes in \emph{coats} and \emph{woolly waistcoats} who \emph{strongly} objected to fresh air! Men are such \textsc{hot-house plants} now-a-days, are they not, quite unlike my dear father, who would never permit a \emph{fire} in the house \emph{before} November the 1st, or \emph{after} March 31st even though the thermometer was at \emph{freezing-point!}

I had \emph{no} difficulty in getting a comfortable room at the Station Hotel, \emph{late} as it was. In the old days, an \emph{unmarried} woman arriving \emph{alone} at \emph{midnight} with a \emph{suitcase} would hardly have been considered \emph{respectable}—what a wonderful difference one finds today! I am \emph{grateful} to have lived to see such changes, because whatever old-fashioned people may say about the greater \emph{decorum} and \emph{modesty} of women in Queen Victoria's time, those who can remember the old conditions know how \emph{difficult} and \emph{humiliating} they were!

Yesterday morning, of course, my \emph{first} object was to find a \emph{suitable boarding-house}, in accordance with your instructions, and I was \emph{fortunate} enough to hit upon this house at the \emph{second} attempt. It is very well run and \emph{refined}, and there are three \emph{elderly ladies} who are \emph{permanent} boarders here, and are \emph{well up} in all the \textsc{gossip} of the town, so that nothing could be more \emph{advantageous} for our purpose!!

As soon as I had engaged my room, I went out for a little \emph{voyage of discovery}. I found a very helpful \emph{policeman} in the High Street, and asked him where to find Mrs~Wrayburn's house. He knew it quite well, and told me to take the \emph{omnibus} and it would be a penny ride to the \buildingname{Fisherman's Arms} and then about 5 minutes' walk. So I followed his directions, and the 'bus took me right into the country to a \emph{cross-roads} with the \buildingname{Fisherman's Arms} at the corner. The conductor was most polite and helpful and showed me the way, so I had \emph{no difficulty} in finding the house.

It is a \emph{beautiful old place}, standing in its own grounds—quite a \emph{big} house built in the \emph{eighteenth century}, with an \emph{Italian} porch and a lovely green lawn with a cedar-tree and formal flower-beds, and in summer must be really a \emph{garden of Eden}. I looked at it from the road for a little time—I did not think this would be at all \emph{peculiar} behaviour, if anybody saw me, because \emph{anybody} might be interested in such a fine old place. Most of the \emph{blinds} were down, as though the greater part of the house were \emph{uninhabited}, and I could not see any \emph{gardener} or anybody about—I suppose there is not very much to be done in the garden this time of the year. One of the \emph{chimneys} was smoking, however, so there were \emph{some} signs of life about the place.

I took a little \emph{walk} down the road and then turned back and passed the house again, and this time I saw a servant just passing round the corner of the house, but of course she was \emph{too far off} for me to speak to. So I took the omnibus back again and had lunch at Hillside View, so as to make acquaintance with my fellow-boarders.

Naturally I did not want to seem \emph{too eager} all at once, so I said nothing about Mrs~Wrayburn's house \emph{at first}, but just talked \emph{generally} about Windle. I had some difficulty in parrying the \emph{questions} of the good ladies, who \emph{wondered} very much \emph{why} a stranger had come to Windle at this time of year, but without telling many actual \emph{untruths} I think I left them with the \emph{impression} that I had come into a little fortune~(!) and was visiting the Lake District to find a suitable spot in which to settle next \emph{summer!} I talked about \emph{sketching}—as girls we were \emph{all} brought up to dabble a little in water-colours, so that I was able to display quite sufficient \emph{technical knowledge to satisfy them!}

That gave me quite a \emph{good} opportunity to ask about the \emph{house!!} Such a \emph{beautiful} old place, I said, and did anybody live there? (Of \emph{course} I did not blurt this out \emph{all at once}—I waited till they had told me of the many \emph{quaint spots} in the district that would interest an artist!) Mrs~Pegler, a very \emph{stout}, \textsc{pussy} old lady, with a \textsc{long tongue}~(!) was able to tell me \emph{all} about it. My dear Lord~Peter, what I do \emph{not} know now about the \emph{abandoned wickedness} of Mrs~Wrayburn's early life is really \textsc{not worth knowing!!} But what was \emph{more to the point} is that she told me the \emph{name of} Mrs~Wrayburn's \emph{nurse-companion}. She is a \textsc{Miss~Booth}, a retired nurse, about \emph{sixty} years old, and she lives \emph{all alone} in the house with Mrs~Wrayburn, except for the \emph{servants}, and a \emph{housekeeper}. When I heard that Mrs~Wrayburn was so \emph{old}, and \emph{paralysed} and \emph{frail}, I said was it not very \emph{dangerous} that Miss~Booth should be the only attendant, but Mrs~Pegler said the housekeeper was a \emph{most trustworthy} woman who had been with Mrs~Wrayburn for many years, and was \emph{quite} capable of looking after her any time when Miss~Booth was out. So it appears that Miss~Booth does go out sometimes! Nobody in this house seems to \emph{know} her \emph{personally}, but they say she is often to be seen in the town in \emph{nurse's uniform}. I managed to extract quite a good description of her, so if I should happen to meet her, I daresay I shall be \emph{smart} enough to \emph{recognise} her!

That is really \emph{all} I have been able to discover in \emph{one} day. I hope you will not be \emph{too} disappointed, but I was obliged to listen to a terrible amount of \emph{local history} of one kind and another, and of course I could not \textsc{force} the conversation round to Mrs~Wrayburn in any suspicious way.

I will let you know as \emph{soon} as I get the \emph{least bit} more information.

\closeletter[Most sincerely yours,]{Katharine Alexandra Climpson.}
\end{mail}

Miss~Climpson finished her letter in the privacy of her bedroom, and secured it carefully in her capacious handbag before going downstairs. A long experience of boarding-house life warned her that to display openly an envelope addressed even to a minor member of the nobility would be to court a quite unnecessary curiosity. True, it would establish her status, but at that moment Miss~Climpson hardly wished to move in the limelight. She crept quietly out at the hall door, and turned her steps towards the centre of the town.

On the previous day, she had marked down one principal tea-shop, two rising and competitive tea-shops, one slightly passé and declining tea-shop, a Lyons, and four obscure and, on the whole, negligible tea-shops which combined the service of refreshments with a trade in sweets. It was now half-past ten. In the next hour and a half she could, with a little exertion, pass in review all that part of the Windle population which indulged in morning coffee.

She posted her letter and then debated with herself where to begin. On the whole, she inclined to leave the Lyons for another day. It was an ordinary plain Lyons, without orchestra or soda-fountain. She thought that its clientèle would be chiefly housewives and clerks. Of the other four, the most likely was, perhaps, the \buildingname{Central.} It was fairly large, well-lighted and cheerful and strains of music issued from its doors. Nurses usually like the large, well-lighted and melodious. But the \buildingname{Central} had one drawback. Anyone coming from the direction of Mrs~Wrayburn's house would have to pass all the others to get to it. This fact unfitted it for an observation post. From this point of view, the advantage lay with \buildingname{Ye Cosye Corner,} which commanded the 'bus-stop. Accordingly, Miss~Climpson decided to start her campaign from that spot. She selected a table in the window, ordered a cup of coffee and a plate of digestive biscuits and entered upon her vigil.

After half an hour, during which no woman in nurse's costume had been sighted, she ordered another cup of coffee and some pastries. A number of people—mostly women—dropped in, but none of them could by any possibility be identified with Miss~Booth. At half-past eleven, Miss~Climpson felt to stay any longer would be conspicuous and might annoy the management. She paid her bill and departed.

The \buildingname{Central} had rather more people in it than \buildingname{Ye Cosye Corner,} and was in some ways an improvement, having comfortable wicker chairs instead of fumed oak settles, and brisk waitresses instead of languid semi-gentlewomen in art-linen. Miss~Climpson ordered another cup of coffee and a roll and butter. There was no window-table vacant, but she found one close to the orchestra from which she could survey the whole room. A fluttering dark-blue veil at the door made her heart beat, but it proved to belong to a lusty young person with two youngsters and a perambulator, and hope withdrew once more. By twelve o'clock, Miss~Climpson decided that she had drawn blank at the \buildingname{Central.}

Her last visit was to the \buildingname{Oriental}—an establishment singularly ill-adapted for espionage. It consisted of three very small rooms of irregular shape, dimly lit by forty-watt bulbs in Japanese shades, and further shrouded by bead curtains and draperies. Miss~Climpson, in her inquisitive way, wandered into all its nooks and corners, disturbing several courting couples, before returning to a table near the door and sitting down to consume her fourth cup of coffee. Half-past twelve came, but no Miss~Booth. <She can't come now,> thought Miss~Climpson, <she will have to get back and give her patient lunch.>

She returned to Hillside View with but little appetite for the joint of roast mutton.

At half-past three she sallied out again, to indulge in an orgy of teas. This time she included the Lyons and the fourth tea-shop, beginning at the far end of the town and working her way back to the 'bus-stop. It was while she was struggling with her fifth meal, in the window of \buildingname{Ye Cosye Corner,} that a hurrying figure on the pavement caught her eye. The winter evening had closed in, and the street-lights were not very brilliant, but she distinctly saw a stoutish middle-aged nurse in a black veil and grey cloak pass along on the nearer pavement. By craning her neck, she could see her make a brisk spurt, scramble on the 'bus at the corner and disappear in the direction of the \buildingname{Fisherman's Arms.}

<How vexatious!> said Miss~Climpson, as the vehicle disappeared. <I must have just missed her somewhere. Or perhaps she was having tea in a private house. Well, I'm afraid this is a blank day. And I do feel so \emph{full} of tea!>

It was fortunate that Miss~Climpson had been blest by Heaven with a sound digestion, for the next morning saw a repetition of the performance. It was possible, of course, that Miss~Booth only went out two or three times a week, or that she only went out in the afternoon, but Miss~Climpson was taking no chances. She had at least achieved the certainty that the 'bus-stop was the place to watch. This time she took up her post at \buildingname{Ye Cosye Corner} at 11 o'clock and waited till twelve. Nothing happened and she went home.

In the afternoon she was there again at three. By this time the waitress had got to know her, and betrayed a certain amused and tolerant interest in her comings and goings. Miss~Climpson explained that she liked so much to watch the people pass, and spoke a few words in praise of the Café and its service. She admired a quaint old inn on the opposite side of the street, and said she thought of making a sketch of it.

<Oh, yes,> said the girl, <there's a many artists comes here for that.>

This gave Miss~Climpson a bright idea, and the next morning she brought a pencil and sketch-book with her.

By the extraordinary perversity of things in general, she had no sooner ordered her coffee, opened the sketch-book and started to outline the gables of the inn, than a 'bus drew up, and out of it stepped the stout nurse in the black and grey uniform. She did not enter \buildingname{Ye Cosye Corner,} but marched on at a brisk pace down the opposite side of the street, her veil flapping like a flag.

Miss~Climpson uttered a sharp exclamation of annoyance, which drew the waitress's attention.

<How provoking!> said Miss~Climpson. <I have left my rubber behind. I must just run out and buy one.>

She dropped the sketch-book on the table and made for the door.

<I'll cover your coffee up for you, miss,> said the girl, helpfully. <Mr~Bulteel's, down near the <Bear,> is the best stationer's.>

<Thank you, thank you,> said Miss~Climpson, and darted out.

The black veil was still flapping in the distance. Miss~Climpson pursued breathlessly, keeping to the near side of the road. The veil dived into a chemist's shop. Miss~Climpson crossed the road a little behind it and stared into a window full of baby-linen. The veil came out, fluttered undecidedly on the pavement, turned, passed Miss~Climpson and went into a boot-shop.

<If it's shoe-laces, it'll be quick,> thought Miss~Climpson, <but if it's trying-on it may be all morning.> She walked slowly past the door. By good luck a customer was just coming out, and, peering past him, Miss~Climpson just caught a glimpse of the black veil vanishing into the back premises. She pushed the door boldly open. There was a counter for sundries in the front of the shop, and the doorway through which the nurse had vanished was labelled <Ladies' Department.>

While buying a pair of brown silk laces, Miss~Climpson debated with herself. Should she follow and seize this opportunity? Trying on shoes is usually a lengthy business. The subject is marooned for long periods in a chair, while the assistant climbs ladders and collects piles of cardboard boxes. It is also comparatively easy to enter into conversation with a person who is trying on shoes. But there is a snag in it. To give colour to your presence in the Fitting department, you must yourself try on shoes. What happens? The assistant first disables you by snatching off your right-hand shoe, and then disappears. And supposing, meanwhile, your quarry completes her purchase and walks out? Are you to follow, hopping madly on one foot? Are you to arouse suspicion by hurriedly replacing your own footgear and rushing out with laces flying and an unconvincing murmur about a forgotten engagement? Still worse, suppose you are in an amphibious condition, wearing one shoe of your own and one of the establishment's? What impression will you make by suddenly bolting with goods to which you are not entitled? Will not the pursuer very quickly become the pursued?

Having weighed this problem in her mind, Miss~Climpson paid for her shoe-laces and retired. She had already bilked a tea-shop, and one misdemeanour in a morning was about as much as she could hope to get away with.

The male detective, particularly when dressed as a workman, an errand-boy or a telegraph-messenger, is favourably placed for <shadowing.> He can loaf without attracting attention. The female detective must not loaf. On the other hand, she can stare into shop-windows for ever. Miss~Climpson selected a hat-shop. She examined all the hats in both windows attentively, coming back to gaze in a purposeful manner at an extremely elegant model with an eye-veil and a pair of excrescences like rabbits'-ears. Just at the moment when any observer might have thought that she had at last made up her mind to go in and ask the price, the nurse came out of the boot-shop. Miss~Climpson shook her head regretfully at the rabbits'-ears, darted back to the other window, looked, hovered, hesitated—and tore herself away.

The nurse was now about thirty yards ahead, moving well, with the air of a horse that sights his stable. She crossed the street again, looked into a window piled with coloured wools, thought better of it, passed on, and turned in at the door of the Oriental Café.

Miss~Climpson was in the position of one who, after prolonged pursuit, has clapped a tumbler over a moth. For the moment the creature is safe and the pursuer takes breath. The problem now is to extract the moth without damage.

It is easy, of course, to follow a person into a café and sit down at her table, if there is room there. But she may not welcome you. She may feel it perverse in you to thrust yourself upon her when other tables are standing empty. It is better to offer some excuse, such as restoring a dropped handkerchief or drawing attention to an open handbag. If the person will not provide you with an excuse, the next best thing is to manufacture one.

The stationer's shop was only a few doors off. Miss~Climpson went in and purchased an indiarubber, three picture post-cards, a BB pencil and a calendar, and waited while they were made up into a parcel. Then she slowly made her way across the street and turned into the \buildingname{Oriental.}

In the first room she found two women and a small boy occupying one recess, an aged gentleman drinking milk in another, and a couple of girls consuming coffee and cakes in a third.

<Excuse me,> said Miss~Climpson to the two women, <but does this parcel belong to you? I picked it up just outside the door.>

The elder woman, who had evidently been shopping, hastily passed in review a quantity of miscellaneous packages, pinching each one by way of refreshing her memory as to the contents.

<I don't think it's mine, but really I can't say for certain. Let me see. That's eggs and that's bacon and—what's this, Gertie? Is that the mouse-trap? No, wait a minute, that's cough-mixture, that is—and that's Aunt Edith's cork soles, and that's Nugget—no, bloater paste, this here's the Nugget—why, bless my soul, I believe I \emph{have} been and gone and dropped the mouse-trap—but that don't look like it to me.>

<No, Mother,> said the younger woman, <don't you remember, they were sending round the mouse-trap with the bath.>

<Of course, so they were. Well, that accounts for that. The mouse-trap and the two frying-pans, they was all to go with the bath, and that's all except the soap, which you've got, Gertie. No, thank you very much, all the same, but it isn't ours; somebody else must have dropped it.>

The old gentleman repudiated it firmly, but politely, and the two girls merely giggled at it. Miss~Climpson passed on. Two young women with their attendant young men duly thanked her in the second room, but said the parcel was not theirs.

Miss~Climpson passed into the third room. In one corner was a rather talkative party of people with an Airedale, and at the back, in the most obscure and retired of all the Oriental nooks and corners, sat the nurse, reading a book.

The talkative party had nothing to say to the parcel, and Miss~Climpson, with her heart beating fast, bore down upon the nurse.

<Excuse me,> she said, smiling graciously, <but I think this little parcel must be yours. I picked it up just in the doorway and I've asked all the other people in the café.>

The nurse looked up. She was a grey-haired, elderly woman, with those curious large blue eyes which disconcert the beholder by their intense gaze, and are usually an index of some emotional instability. She smiled at Miss~Climpson and said pleasantly:

<No, no, it isn't mine. So kind of you. But I have all my parcels here.>

She vaguely indicated the cushioned seat which ran round three sides of the alcove, and Miss~Climpson, accepting the gesture as an invitation, promptly sat down.

<How very odd,> said Miss~Climpson, <I made sure someone must have dropped it coming in here. I wonder what I had better do with it.> She pinched it gently. <I shouldn't think it was valuable, but one never knows. I suppose I ought to take it to the police-station.>

<You could hand it to the cashier,> suggested the nurse, <in case the owner came back here to claim it.>

<Well now, so I could,> cried Miss~Climpson. <How clever of you to think of it. Of course, yes, that would be the best way. You must think me very foolish, but the idea never occurred to me. I'm not a very practical person, I'm afraid, but I do so admire the people who are. I should never do to take up \emph{your} profession, should I? Any little emergency leaves me \emph{quite} bewildered.>

The nurse smiled again.

<It is largely a question of training,> she said. <And of \emph{self}-training, too, of course. All these little weaknesses can be cured by placing the mind under a Higher Control—don't you believe that?>

Her eyes rested hypnotically upon Miss~Climpson's.

<I suppose that is true.>

<It is such a mistake,> pursued the nurse, closing her book and laying it down on the table, <to imagine that anything in the mental sphere is large or small. Our least thoughts and actions are equally directed by the higher centres of spiritual power, if we can bring ourselves to believe it.>

A waitress arrived to take Miss~Climpson's order.

<Oh, dear! I seem to have intruded myself upon your table\longdots>

<Oh, don't get up,> said the nurse.

<Are you sure? Really? because I don't want to interrupt you\longdash>

<Not at all. I live a very solitary life, and I am always glad to find a friend to talk to.>

<How nice of you. I'll have scones and butter, please, and a pot of tea. This is such a nice little café, don't you think?—so quiet and peaceful. If only those people wouldn't make such a noise with that dog of theirs. I don't like those great big animals, and I think they're quite dangerous, don't you?>

The reply was lost on Miss~Climpson, for she had suddenly seen the title of the book on the table, and the Devil, or a Ministering Angel (she was not quite sure which) was, so to speak, handing her a fullblown temptation on a silver salver. The book was published by the Spiritualist Press and was called <\workname{Can the Dead Speak?}>

In a single moment of illumination, Miss~Climpson saw her plan complete and perfect in every detail. It involved a course of deception from which her conscience shrank appalled, but it was certain. She wrestled with the demon. Even in a righteous cause, could anything so wicked be justified?

She breathed what she thought was a prayer for guidance, but the only answer was a small whisper in her ear, <Oh, jolly good work, Miss~Climpson!> and the voice was the voice of Peter Wimsey.

<Pardon me,> said Miss~Climpson, <but I see you are a student of spiritualism. How interesting that is!>

If there was one subject in the world about which Miss~Climpson might claim to know something, it was Spiritualism. It is a flower which flourishes bravely in a boarding-house atmosphere. Time and again, Miss~Climpson had listened while the apparatus of planes and controls, correspondences and veridical communications, astral bodies, auras and ectoplastic materialisations was displayed before her protesting intelligence. That to the Church it was a forbidden subject she knew well enough, but she had been paid companion to so many old ladies and had been forced so many times to bow down in the House of Rimmon.

And then there had been the quaint little man from the Psychical Research Society. He had stayed a fortnight in the same private hotel with her at Bournemouth. He was skilled in the investigation of haunted houses and the detection of poltergeists. He had rather liked Miss~Climpson, and she had passed several interesting evenings hearing about the tricks of mediums. Under his guidance she had learnt to turn tables and produce explosive cracking noises; she knew how to examine a pair of sealed slates for the marks of the wedges which let the chalk go in on a long black wire to write spirit-passages. She had seen the ingenious rubber gloves which leave the impression of spirit hands in a bucket of paraffin-wax, and which, when deflated, can be drawn delicately from the hardened wax through a hole narrower than a child's wrist. She even knew theoretically, though she had never tried it, how to hold her hands to be tied behind her back so as to force that first deceptive knot which makes all subsequent knots useless, and how to flit about the room banging tambourines in the twilight in spite of having been tied up in a black cabinet with both fists filled with flour. Miss~Climpson had wondered greatly at the folly and wickedness of mankind.

The nurse went on talking, and Miss~Climpson answered mechanically.

<She's only a beginner,> said Miss~Climpson to herself. <She's reading a text-book\dots And she is quite uncritical\dots Surely she knows that that woman was exposed long ago\dots People like her shouldn't be allowed out alone—they're living incitements to fraud\dots I don't know this Mrs~Craig she is talking about, but I should say she was as twisty as a corkscrew\dots I must avoid Mrs~Craig, she probably knows too much \dots if the poor deluded creature will swallow that, she'll swallow anything.>

<It does seem \emph{most} wonderful, doesn't it?< said Miss~Climpson, aloud. >But isn't it a wee bit \emph{dangerous}? I've been told I'm sensitive myself, but I have never dared to \emph{try}. Is it \emph{wise} to open one's mind to these supernatural influences?>

<It's not dangerous if you know the right way,> said the nurse. <One must learn to build up a shell of pure thoughts about the soul, so that no evil influences can enter it. I have had the most marvellous talks with the dear ones who have passed over\longdots>

Miss~Climpson refilled the tea-pot and sent the waitress for a plate of sugary cakes.

<\dots unfortunately I am not mediumistic myself—not yet, that is. I can't get anything when I'm alone. Mrs~Craig says that it will come by practice and concentration. Last night I was trying with the Ouija board, but it would only write spirals.>

<Your conscious mind is too active, I expect,> said Miss~Climpson.

<Yes, I daresay that is it. Mrs~Craig says that I am wonderfully sympathetic. We get the most wonderful results when we sit together. Unfortunately she is abroad just now.>

Miss~Climpson's heart gave a great leap, so that she nearly spilled her tea.

<You yourself are a medium, then?> went on the nurse.

<I have been told so,> said Miss~Climpson, guardedly.

<I wonder,> said the nurse, <whether if we sat together\longdash>

She looked hungrily at Miss~Climpson.

<I don't really like\longdash>

<Oh, do! You are such a sympathetic person. I'm sure we should get good results. And the spirits are so pathetically anxious to communicate. Of course, I wouldn't like to try unless I was sure of the person. There are so many fraudulent mediums about>—(<So you do know that much!> thought Miss~Climpson)—<but with somebody like yourself one is absolutely safe. You would find it made such a difference in your life. I used to be so unhappy over all the pain and misery in the world—we see so much of it, you know—till I realised the certainty of survival and how all our trials are merely sent to fit us for life on a higher plane.>

<Well,> said Miss~Climpson, slowly, <I'm willing just to try. But I can't say I really \emph{believe} in it, you know.>

<You would—you would.>

<Of course, I've seen one or two strange things happen—things that couldn't be tricks, because I knew the people—and which I couldn't explain\longdash>

<Come up and see me this evening, now do!> said the nurse, persuasively. <We'll just have one quiet sitting and then we shall see whether you really are a medium. I've no doubt you are.>

<Very well,> said Miss~Climpson. <What is your name, by the way?>

<Caroline Booth—Miss~Caroline Booth. I'm nurse to an old, paralysed lady in the big house along the Kendal Road.>

<Thank goodness for that, anyway,> thought Miss~Climpson. Aloud she said:

<And my name is Climpson; I think I've got a card somewhere. No—I've left it behind. But I'm staying at Hillside View. How do I get to you?>

Miss~Booth mentioned the address and the time of the 'bus, and added an invitation to supper, which was accepted. Miss~Climpson went home and wrote a hurried note:

\begin{mail}{}{My dear Lord~Peter—}

I am sure you have been \emph{wondering} what has \emph{happened} to me. But \emph{at last} I have \textsc{news!} I have \textsc{stormed the citadel!!!} I am going to the \emph{house tonight} and you may expect \textsc{great things!!!}

In \emph{haste},

\closeletter[Yours very sincerely,]{Katharine A. Climpson.}
\end{mail}

Miss~Climpson went out into the town again after lunch. First, being an honest woman, she retrieved her sketch-book from \buildingname{Ye Cosye Corner} and paid her bill, explaining that she had run across a friend that morning and been detained. She then visited a number of shops. Eventually she selected a small metal soap-box which suited her requirements. Its sides were slightly convex, and when closed and pinched slightly, it sprang back with a hearty cracking noise. This, with a little contrivance and some powerful sticking-plaster, she fixed to a strong elastic garter. When clasped about Miss~Climpson's bony knee and squeezed sharply against the other knee, the box emitted a series of cracks so satisfying as to convince the most sceptical. Miss~Climpson, seated before the looking-glass, indulged in an hour's practice before tea, till the crack could be produced with the minimum of physical jerk.

Another purchase was a length of stiff black-bound wire, such as is used for making hat-brims. Used double, neatly bent to a double angle and strapped to the wrist, this contrivance was sufficient to rock a light table. The weight of a heavy table would be too much for it, she feared, but she had had no time to order blacksmith's work. She could try, anyway. She hunted out a black velvet rest-gown with long, wide sleeves, and satisfied herself that the wires could be sufficiently hidden.

At six o'clock, she put on this garment, fastened the soap-box to her leg—turning the box outward, lest untimely cracks should startle her fellow-travellers, muffled herself in a heavy rain-cloak of Inverness cut, took hat and umbrella and started on her way to steal Mrs~Wrayburn's will.